%% Chapter 8 --- Conclusion (English version; labels match the Portuguese deposit version).
\chapter{Conclusion}
\label{chap:conclusao}

This chapter answers the three research questions with the evidence
of Chapters~\ref{chap:benchmark} to \ref{chap:destilacao}, enumerates
the contributions --- with the negative results on the same footing
as the positive ones --- translates the results into guidance for the
practitioner, consolidates the limitations and threats to validity of
the study as a whole, and points to future work.

\section{Final remarks}
\label{sec:conc-finais}

This dissertation started from the question that organized a decade
of literature on tabular data (``do trees still win?'') and ended
answering a different one, because the evidence collected showed that
the original had expired. The three research questions stated in
Chapter~\ref{chap:introducao} receive here their direct answers, each
with the evidence of the chapter that answered it.

\textbf{RQ1} --- \emph{does the accuracy tie among the families of
tabular models survive a rigorous protocol and the arrival of the
2026 generation of foundation models?}
Chapter~\ref{chap:benchmark} answers in two beats. Within the
pre-2026 generation, yes: the top is practical equivalence gauged
under the declared criteria (8/10 equivalent under a ROPE of
$\pm0.01$ on the lenient $E[\theta]>0.5$ criterion, and 3/10 --- the
GBDTs --- at the strict 0.95 cutoff; sensitivity
$0.005/0.01/0.02 \to 2/8/10$ on the lenient criterion), with the
model$\times$dataset interaction dominating the rank variance (0.69,
against 0.20 for the family and 0.11 for the model-within-family), so
that the ``GBDT vs.\ deep learning'' axis is a weak abstraction for
these data. Across generations, no: the pre-trained paradigm pulls
away from the pack (TabFM zero-shot with average ranks 2.7/1.0/1.2 in
the three tasks --- multiclass over the two datasets where it
operates --- and best model on 13 of the 18 datasets; TabPFN as the
consistent runner-up). That lead is supported by the consistency
across datasets and the paired effect sizes --- not by the
underpowered frequentist post-hoc nor by the Bayesian instrument,
which does not separate TabFM from XGBoost and CatBoost in practical
equivalence --- and comes at the price of four to five orders of
magnitude in serving latency (43{,}262 against XGBoost's
0.33~\us/row under the standardized protocol). The instance layer (Section~\ref{sec:cap5-ato5}) further
says \emph{where} the shift lives: XGBoost and the foundation models
agree on ${\ge}93.6\%$ of instances (median), the disputed fringe is
converted in a strongly heterogeneous way across datasets, and the
easy explanations (low margin, minority class, little data) were not
confirmed under the pre-registered criteria.

\textbf{RQ2} --- \emph{which decision artifact can a benchmark with a
realistic $N$ of datasets actually support, a learned router or a
constraint-based policy --- and how does that answer change when the
decision weighs accuracy against cost?} A constraint-based policy,
not a learned router. Under leave-one-dataset-out with 18 datasets,
the meta-feature router scores below the trivial majority-class rule
(0.11 against 0.44 with 11 models; 0.72 against 0.78 with 14), even
though the conditional reversals it should exploit are real ---
recent deep learning takes the multiclass top (descriptive, $N{=}3$),
the GBDTs gain advantage under severe imbalance, TabPFN dominates the
low-data regime ($n \leq 4{,}000$; and the within-dataset subsampling
shows that the foundation models' advantage is not confined to small
datasets, Section~\ref{sec:cap5-ato5}). What the data support is a
framework of constraints verifiable before any training (latency,
architectural limits, interpretability), with the ``foundation model
first'' policy as the zero-median-regret default on the primary
metric (a descriptive evaluation, on the same 17 datasets). When the decision weighs accuracy against cost, the answer
changes early: the recommendation stops being the foundation model
between $w{=}0.9$ and $w{=}0.8$ of accuracy weight, under the rank
normalization adopted (experiment E2,
Chapter~\ref{chap:framework}) --- and that is why the flowchart's
first question is the latency budget: it is the one that changes the
answer earliest.

\textbf{RQ3} --- \emph{what does knowledge distillation transfer from
a tabular foundation model in regression, at what serving cost, and
under which conditions verifiable a priori?} It transfers part of the
shape of the predictive distribution, not the mean --- with
TabPFN-2.5 as the sole distributional teacher. The essential of what
the teacher knows about the mean, a tuned GBDT learns on its own from
the raw labels: the student trained on the teacher's means beats the
control on 4 of 20 datasets with TabPFN, and the replication with
TabFM as teacher, on the five central datasets, remains negative. The
one trained on the quantile curves retains, at the median, 19\% of
the CRPS gap with inherited hyperparameters (12\% with both arms
re-tuned) --- a median CRPS gain of 2.8\% shrinking to 1.4\% --- with
positive retention on 16 of 20 datasets (sign test,
$p\approx0.019$ under the two-stage design's sequential correction),
and beats the best of QRF and NGBoost on 17 of 20 (experiment E1).
The distributional student's serving cost is on the order of
$10^{2}$~\us{} per row in the measured frontier cases (99 and
169~\us/row; Section~\ref{sec:destil-fronteira}) --- above the point
control, and $2.5\times$ to $27\times$ below the teacher's cheapest
competitive configuration where the teacher's absolute advantage is
large. The
verifiable conditions are three, and the out-of-fold labeling that
measures them costs minutes: transform the target, test a strong
native distributional baseline, and distill only the advantage that
remains to the teacher --- with OOF labeling itself as the student's
calibration safeguard (0.5--5.1~p.p.\ of PICP80,
Chapter~\ref{chap:destilacao}).

\section{Contributions}
\label{sec:conc-contribuicoes}

The five contributions stated in
Section~\ref{sec:intro-contribuicoes} are revisited here in the same
order; under each are listed the results that compose it, with the
nature of each --- \emph{evidence} (an empirical result under a
declared protocol), \emph{framing} (a verified gap delimitation), or
\emph{rule} (an operational prescription derived from the evidence)
--- and the chapter of origin. Negative results are listed with the
same status as positive ones, per the study's conduct contract
(Chapter~\ref{chap:metodologia}).

\begin{enumerate}
  \item \textbf{The multi-criteria benchmark of the 2026 generation.}
  \begin{enumerate}
  \item \textbf{Neutral 14-model benchmark (evidence;
    Chapter~\ref{chap:benchmark}).} 14 models $\times$ 18 datasets
    (250 of the 252 cells; two structural exclusions), with a
    hold-out untouched by selection, frequentist \emph{and} Bayesian
    inference with ROPE and sensitivity analysis, and three cost
    axes under the same protocol. From it come the previous
    generation's positive tie, the shift of the top by the 2026
    generation (a consistent lead, not a statistically confirmed
    separation) and, as far as the dated checks reach, the first
    neutral TabPFN~$\times$~TabFM comparison under an identical
    protocol, framed as an inter-laboratory replication of the
    paradigm, with pinned versions and mirrored policies.
  \item \textbf{Instance-level complementarity analysis (evidence;
    Chapter~\ref{chap:benchmark}, Section~\ref{sec:cap5-ato5}).}
    Pre-registered with dated falsification criteria and exact
    reproduction of the benchmark's final models (52/52 cells
    through a sanity gate), it establishes that XGBoost and the
    foundation models agree on ${\ge}93.6\%$ of instances (median),
    that the FM's advantage is strongly heterogeneous across
    datasets (exact McNemar per pair) and, as informative negatives,
    it does not confirm the explanations of low margin, minority
    class, and advantage confined to little data (the last with a
    contrary descriptive sign in the aggregate).
  \end{enumerate}
  \item \textbf{The independent test of the Kolmogorov--Arnold
    family.}
  \begin{enumerate}
  \item \textbf{Unfavorable result under a uniform protocol
    (negative evidence; Chapter~\ref{chap:benchmark}).} In the
    family's first confrontation with the full
    GBDT/DL/foundation-model grid under a multi-criteria protocol
    --- as far as the dated checks reach --- KAN and TabKAN finish
    in the bottom third on all three tasks (ranks 11.0/12.0/9.8 and
    11.4/13.3/11.8), in line with the cost-parity critique
    \citep{kancritical2024} and against the proponents' claims ---
    and against \citet{poeta2024}'s KAN-favorable result, obtained
    without GBDT baselines and under short training, which the
    uniform protocol does not reproduce; the recorded counterpart is
    that both are cheap to serve (2.3--3.1~\us/row) --- the problem
    is performance, not cost.
  \end{enumerate}
  \item \textbf{The evaluated decision framework.}
  \begin{enumerate}
  \item \textbf{Negative result of meta-feature routing (negative
    evidence; Chapter~\ref{chap:framework}).} Under LODO with
    $N{=}18$, a learned router does not beat the majority-class
    baseline in either generation (0.11 vs.\ 0.44 with 11 models;
    0.72 vs.\ 0.78 with 14) --- corroborated by independent
    contemporary evidence \citep{herre2026}. The design consequence:
    the framework contains no winning-family predictor.
  \item \textbf{Decision framework evaluated by regret
    (evidence~$+$~rule; Chapter~\ref{chap:framework}).} The fixed
    policy ``foundation model first'' is regret-optimal on the
    primary metric (a descriptive evaluation: median regret 0.000 in
    the 2026 generation,
    preserved in the executable always-TabFM policy --- mean 0.080,
    worst case 1.000 on \texttt{telco}, the latter largely an
    artifact of the preprocessing measured in
    Chapter~\ref{chap:metodologia}), and the final framework is a
    flowchart of four constraints verifiable \emph{a priori} ---
    the form the data support.
  \item \textbf{Weighted multi-criteria regret --- E2 (evidence;
    Chapter~\ref{chap:framework}).} The registered experiment that
    locates the recommendation's inversion point: ``FM first'' is
    optimal in median only for $w \geq 0.9$ of accuracy weight;
    with $w \leq 0.6$, ``always-GBDT'' has a zero median regret. It
    is the quantitative justification of the flowchart's
    architecture --- latency is the first question because it is
    the one that changes the answer earliest.
  \end{enumerate}
  \item \textbf{The distributional distillation of tabular
    foundation models.}
  \begin{enumerate}
  \item \textbf{The framing and the central result
    (framing~$+$~evidence; Chapter~\ref{chap:destilacao}).} The
    framing --- transferring quantile curves \emph{from a tabular
    foundation-model teacher} to a GBDT student, evaluating by CRPS
    and coverage --- was verified as having no prior
    \emph{academic} work in this study's dated searches, July to
    November 2026 (complete record in
    Appendix~\ref{app:suplementares}; a public repository of
    September 2026 --- referenced in
    Section~\ref{sec:rel-destilacao} --- distills TabFM into
    XGBoost for classification and \emph{point} regression, without
    touching the distributional framing --- the non-tabular lineage
    of distributional distillation is cited as a precedent in
    Section~\ref{sec:rel-destilacao}). The evidence: median
    retention of 12\% to 19\% of the CRPS gap, depending on how the
    arms are tuned, positive on 16 of 20 datasets ($p=0.006$,
    $p\approx0.019$ under the sequential correction; Bayesian
    posterior probability of 0.96 on positive retention, with
    expected posterior mass 0.70), including a \emph{born-again}
    student (which matches its own teacher to within $10^{-4}$ of
    CRPS), serving at about a hundred microseconds per row.
  \item \textbf{Refutation of point distillation at realistic pool
    scale (negative evidence; Chapter~\ref{chap:destilacao}).} The
    distilled point student beats the control on only 4 of 20
    datasets with TabPFN; the replication with TabFM --- whose
    point advantage over the baseline is nearly $6\times$ larger on
    the measured dataset (the binarized-target
    \texttt{california\_housing}) --- remains negative on the five
    central datasets, and the pool-size sweep isolates the
    moderator (residual gain confined to $n \leq 800$,
    dataset-dependent) --- a reconciliation hypothesis with the
    contemporary point gains of \citet{jeong2026}, whose recipe
    differs from ours on axes we did not replicate (full-context
    labeling, synthetic queries, a student trained only on teacher
    predictions).
  \item \textbf{Three-strategy serving frontier (evidence;
    Chapter~\ref{chap:destilacao}).} Distilling, compressing the
    context and shrinking the ensemble measured for the first time
    on the same hold-outs: where the teacher's absolute advantage is
    large, the distilled student is, among the teacher-derived
    strategies,
    the CRPS-optimal distributional system below its cheapest
    competitive configuration ($2.5\times$ and $27\times$ the
    student's latency in the two measured cases --- E1 shows a
    native QRF also inhabits the regime); where a single-member
    teacher already reaches the student's tier, shrinking the
    ensemble suffices. Two engineering side findings: one TabPFN
    member $\approx$ eight at a seventh of the latency, and the
    last context doubling is disproportionately expensive on 16~GB
    GPUs.
  \item \textbf{Deconfounded ablation of out-of-fold labeling
    (evidence and rule; Chapter~\ref{chap:destilacao}).} In
    regression, teacher leakage does not degrade accuracy --- it
    degrades \emph{calibration} (a loss of 0.5 to 5.1~p.p.\ of
    PICP80 at identical context): a deconfounded refinement of
    Pocket~FM's classification finding \citep{pocketfm2026}.
  \item \textbf{Confrontation with classical distributional
    natives --- E1 (evidence; Chapter~\ref{chap:destilacao}).} The
    distilled student beats the best of QRF and NGBoost on 17 of 20
    datasets ($p=0.0013$, one-sided), with the three upsets
    diagnosable --- on \texttt{wine\_quality}, where the
    pre-registered gate already advised against distilling, the
    default QRF beats even the teacher. The negative complement has
    value of its own: the ``test a strong native first'' step of
    the decision rule stopped being rhetorical.
  \item \textbf{Architectural negative of TabFM as a distributional
    teacher (negative evidence; Chapter~\ref{chap:destilacao}).}
    TabFM's regression head is point-valued --- an architectural
    check of July 2026, confirmed by the official technical report
    \citep{kong2026} --- and the spread of its ensemble is
    measurably useless as a predictive distribution (CRPS 0.524
    against TabPFN's 0.069; 16.3\% coverage at a nominal 80\% ---
    measured on one dataset, the binarized-target
    \texttt{california\_housing}$^\dagger$): a
    point leader is not a distributional leader, and the
    distributional head is the decisive architectural asset for
    distillation.
  \item \textbf{Decision rule for distilling (rule;
    Chapter~\ref{chap:destilacao}).} Derived from two hypotheses
    that did not hold (point: refuted; heavy tail: not confirmed)
    and from E1: transform the target, test a strong native, and
    distill only the genuine advantage that survives --- a
    precondition measurable by a minutes-long OOF labeling. It is
    the operational form of the answer to RQ3.
  \end{enumerate}
  \item \textbf{Reusable infrastructure and conduct (transversal;
    Chapters~\ref{chap:metodologia} to~\ref{chap:destilacao}).} An
    end-to-end versioned protocol --- configurations, seeds, data
    hashes, pre-registrations with dated addenda, dated novelty
    checks and public errata --- in which every number in the text
    traces to an artifact. The contribution here is one of
    \emph{process}: the conduct contract determined concrete scope
    decisions, and its correction mechanisms were exercised in real
    cases.
\end{enumerate}

\section{Advice to the practitioner}
\label{sec:conc-conselhos}

The operational distillate of Chapters~\ref{chap:framework}
and~\ref{chap:destilacao}, in the order the questions should be
asked --- each item traces to the result that supports it.

\begin{itemize}
  \item \textbf{Start from the constraints, not from performance.}
    Four questions verifiable before any training: latency budget,
    architectural limits, interpretability requirement, and only
    then the default (flowchart of
    Section~\ref{sec:arcabouco-resultante}). Do not try to learn a
    meta-feature router with a few dozen datasets: under LODO, it
    does not beat the trivial rule
    (Section~\ref{sec:resultado-negativo}).
  \item \textbf{If microsecond latency matters, GBDT.} With 40\% or
    more of the utility weight on latency, ``always-GBDT'' has a
    zero median regret; the ``FM first'' recommendation already
    flips with 10--20\% of the weight on latency (by the median;
    E2, Section~\ref{sec:regret-ponderado}).
  \item \textbf{If no constraint is active, foundation model
    first.} Median regret 0.000 in the 2026 generation
    (Section~\ref{sec:resultado-positivo}) --- with TabPFN as the
    mature option and TabFM as the frontier, at the versions and
    with the maturity caveats declared in
    Chapter~\ref{chap:benchmark}.
  \item \textbf{Check the architectural limits early.} More than 10
    classes excludes TabFM; 100 classes exclude both foundation
    models (the \texttt{helena} case); context overflow is an
    objective restriction --- these are structural exclusions, not
    details (Section~\ref{sec:arcabouco-resultante}).
  \item \textbf{In regression with uncertainty: transform before
    distilling.} On price-scale targets, a one-line logarithmic
    transformation is free and changes the error regime of every
    system --- it does not eliminate the teacher's relative
    advantage, but it is the first test to run, and it costs
    minutes (Section~\ref{sec:destil-regra}).
  \item \textbf{Test the QRF before distilling.} A default
    distributional native costs seconds of training, serves at
    microseconds and, on a predictable minority of datasets,
    settles the problem before any teacher needs to be loaded ---
    on \texttt{wine\_quality}, it beat even the teacher (E1,
    Section~\ref{sec:destil-e1}); a quantile MLP covers the small,
    numeric datasets (Section~\ref{sec:destil-regra}).
  \item \textbf{Distill when the advantage survives the
    transformations and a strong native --- and measure it first.}
    A single OOF labeling, at a cost of minutes, directly measures
    the teacher's genuine advantage; distill only if it persists
    (Section~\ref{sec:destil-regra}). Do not distill the mean: at
    realistic pools, there is nothing to transfer on the point axis
    (Section~\ref{sec:destil-ponto}).
  \item \textbf{Always label out-of-fold.} The cost is modest, and
    the safeguard protects what accuracy metrics do not see: the
    student's interval coverage (0.5--5.1~p.p.\ of PICP80;
    Section~\ref{sec:destil-oof}).
  \item \textbf{Before paying for the full teacher, shrink the
    ensemble.} A single TabPFN member approximated the eight at a
    seventh of the latency --- sometimes at no accuracy cost ---
    and, when a single teacher member already reaches the student's
    latency tier, that lever closes the case without any
    distillation (Section~\ref{sec:destil-fronteira}).
  \item \textbf{On 16~GB GPUs, distrust the last context
    doubling.} The final doubling, from 25{,}000 to 50{,}000 rows,
    cost $3.3\times$ the latency for a marginal accuracy gain
    (Section~\ref{sec:destil-fronteira}); reserve the full-context
    teacher for batch processing.
\end{itemize}

The list above orders the questions; it remains to name the
scenarios. Table~\ref{tab:conc-perfis} consolidates the same results
as \emph{usage profiles}: for each condition --- six of the ten
verifiable before any training, the rest after one cheap, declared
measurement --- the recommendation the data support, the measured
evidence
that supports it and --- the column the usual benchmarks omit ---
what one accepts to lose in exchange. The final column is the point:
in several profiles, the recommendation is deliberately a model
\emph{worse on the primary metric}, and the trade is defensible
because it was measured. Two results support this accounting: in the
core's pack, giving up the top typically costs less than one
percentage point of AUC (Section~\ref{sec:cap5-ato1}), while the
differences in latency, minority-class recall, and tail cost are of
whole factors to orders of magnitude (columns 3 and 4 of the table).
When the metric on which the models differ is the one the
application pays, it is the one that decides --- and the table says
which it is, case by case.

\begin{table}[htbp]
\centering
\caption[Usage profiles: condition, recommendation, evidence, and
trade]{Usage profiles: verifiable condition, recommendation,
evidence, and accepted cost in exchange. Source:
Chapters~\ref{chap:benchmark}--\ref{chap:destilacao}, sections
indicated.}
\label{tab:conc-perfis}
\scriptsize
\setlength{\tabcolsep}{5pt}
\renewcommand{\arraystretch}{1.15}
\begin{tabular}{p{0.21\linewidth}p{0.17\linewidth}p{0.33\linewidth}p{0.19\linewidth}}
\toprule
\textbf{Profile (condition)} & \textbf{Recommendation} &
\textbf{Evidence} & \textbf{Accepted in exchange} \\
\midrule
No active constraint; latency-tolerant &
Foundation model first (TabPFN mature; TabFM frontier) &
Median regret 0.000 with 14 models in
Chapter~\ref{chap:framework}'s evaluation (the family's ceiling,
measured on the same 18 datasets; the executable always-TabFM
preserves the median, with mean 0.080 and worst case 1.0;
always-TabPFN: median 0.242); absolute best on 13 of 18
(Sections~\ref{sec:resultado-positivo}, \ref{sec:cap5-ato4}) &
Serving cost (from ${\sim}3{,}900\times$ the slowest GBDT ---
TabPFN over LightGBM, Table~\ref{tab:app-custos} --- to
$131{,}097\times$ XGBoost); TabPFN-2.5's non-commercial license \\
\addlinespace
Little data ($n \lesssim 4{,}000$) &
TabPFN &
Average rank 1.0 in the regime, against XGBoost's 2.8, on the 5
swept datasets (Section~\ref{sec:cap5-ato3}) &
Inference latency; large-$n$ leadership depends on the dataset \\
\addlinespace
Rare class with expensive false negatives &
FM at a fixed threshold --- \emph{or} re-optimize the GBDT's
threshold; measure recall at equalized precision before switching &
\texttt{amazon\_employee}: accuracies 0.943--0.949
indistinguishable; minority recall 0.032 $\to$ 0.298 at threshold
0.5 (Section~\ref{sec:cap5-ato5}). Counterweights: the
rare-class-gain hypothesis was not confirmed; the dataset's best AUC
belongs to a GBDT (CatBoost 0.887, Table~\ref{tab:app-binaria}) ---
a fixed-threshold advantage, on one dataset &
Precision (0.706 $\to$ 0.621) and latency; the advantage is
frequent, not law (\texttt{give\_me\_some\_credit}: recalls that
differ by less than 0.02 from one another) \\
\addlinespace
Regression with tail decisions and a \us{} budget &
Distilled student (conformal coverage correction: piloted on 5
datasets) &
Positive retention on 16/20 ($p{=}0.006$ single-stage; $0.019$
sequential), held
re-tuned; at $\tau{=}0.90$, 18/20 ($p{=}0.0002$), ${\sim}8\%$ gain
in the tails; the distilled$+$conformal combination not measured
jointly on the 20 (Sections~\ref{sec:destil-dist},
\ref{sec:destil-e3}) &
Nothing on the point axis (4 of 20); latency above the control, on
the order of $10^2$~\us{} per row; the OOF pipeline; the teacher's
non-commercial license reaches the distilled student \\
\addlinespace
The application's metric inverts the winner &
Decide by the metric the application pays; report both layers &
\texttt{telco}: TabFM loses accuracy (0.771 vs.\ 0.801) and the
McNemar fringe ($-29.8$/thousand), but wins recall (0.663 vs.\
0.570) and balanced accuracy (0.737 vs.\ 0.727)
(Section~\ref{sec:cap5-ato5}) &
The primary metric's ``winner''; an unoptimized threshold; and
\texttt{TotalCharges} (Ch.~4) reaches this dataset: with the column
numeric, TabFM jumps from last to best AUC among the 14 re-evaluated \\
\addlinespace
Advantage heterogeneous per dataset (instance-level disagreement
fringe, Section~\ref{sec:cap5-ato5}) &
Choose per dataset/task, never per instance; no learned router &
Agreement $\geq$0.936; TabFM converts the fringe on \texttt{jannis}
($+59.6$/thousand) and loses it on \texttt{telco}; the LODO router
below the trivial rule (Sections~\ref{sec:cap5-ato5},
\ref{sec:resultado-negativo}) &
The promise of ``a universal winner''; the advantage does not
shrink with size \\
\addlinespace
Microseconds-per-row latency required &
Native GBDT; distillation recovers part of the distributional
advantage &
0.33~\us/row (XGBoost) vs.\ 43{,}262 (TabFM), $131{,}097\times$;
the FM recommendation flips with 10--20\% of latency weight (by
median regret; E2)
(Sections~\ref{sec:cap5-ato2}, \ref{sec:regret-ponderado}) &
The top rank --- a loss $<$1~p.p.\ of AUC in the core's pack (ROPE
vs.\ TabPFN: 8/10 lenient, 3/10 strict); vs.\ TabFM the equivalence
is uncertain ($P$ 0.33; $E[\theta]$ 0.40) \\
\addlinespace
Intrinsic interpretability required (regulation) &
GBDT &
The only family strong in latency \emph{and} interpretability in
the multi-criteria matrix --- a qualitative per-family attribution,
declared, not measured; gate 3 of the flowchart
(Section~\ref{sec:arcabouco-resultante}) &
The top rank, under the same ROPE accounting; no distillation (the
student inherits labels from an opaque teacher) \\
\addlinespace
Architectural limit ($>$10 classes; context overflow) &
GBDT or recent DL depending on the task &
\texttt{helena} (100 classes) excludes both FMs; descriptive
multiclass top: STab/TabM 3.67 vs.\ XGBoost 6.0, $N{=}3$
(Sections~\ref{sec:cap5-ato3}, \ref{sec:arcabouco-resultante}) &
Nothing: the exclusion is structural, not a preference \\
\addlinespace
Small teacher advantage; price-scale target; strong native
available &
Do \emph{not} distill: QRF first, $\log(1{+}y)$, ensemble
shrinking &
E1: the QRF beats the distilled student on 3 of 20 (on
\texttt{wine\_quality}, even the teacher); \texttt{kin8nm}: 1
member at the student's latency; the log changes the error regime
for free, without eliminating the teacher's relative advantage
(Sections~\ref{sec:destil-e1}, \ref{sec:destil-fronteira},
\ref{sec:destil-regra}) &
The whole distillation pipeline --- and that is a gain, not a loss \\
\bottomrule
\end{tabular}
\end{table}

Three readings close the table. First, in the profiles where the
recommendation goes against the primary metric (latency,
interpretability, rare class), the trade is accounted for in the two
final columns: the loss on the primary is bounded by measurement (by
the ROPE in the pack, by the declared threshold in the operational
metrics), while the gain on the application's metric is of whole
factors to orders of magnitude. Second, the negative profiles are as
informative as the positive ones: the table says when \emph{not} to
use the foundation model (microseconds, 100 classes, regulation) and
when \emph{not} to distill (\texttt{kin8nm},
\texttt{wine\_\allowbreak quality}), with the cases measured in the
benchmark itself. Third, no row is a ranking: six are conditions
known before any training, in the form Chapter~\ref{chap:framework}
evaluated; the other four (rare class, metric inversion,
instance-level fringe, teacher advantage) call for a cheap, declared
measurement --- a second metric on the hold-out, an out-of-fold
labeling --- before deciding.

\section{Limitations and threats to validity}
\label{sec:conc-limitacoes}

Each chapter declared its local limitations; this section
consolidates those of the study as a whole, organized by the three
canonical validity types (Section~\ref{sec:met-limitacoes}). Two
transversal limitations enter the taxonomy with their own
prominence: multiplicity and the pure distilled student's
undercoverage, below, under internal validity; the teachers'
licenses, under external validity.

\subsection{Internal validity}

\textbf{Benchmark-aware priors --- the unaddressable residual
risk.} TabFM declares exclusively synthetic pre-training; the
default TabPFN v2.5 classifier checkpoint, however, is fine-tuned
on 43 real datasets that include this benchmark's \texttt{higgs}
--- and, by a name match not verified row by row, probably
\texttt{cardiovascular\_disease} --- (with the sensitivity measured
on the twelve classification datasets: AUC differences between
$-0.0051$ and $+0.0002$, and on \texttt{higgs} the synthetic-only
checkpoint is even slightly better, 0.8113 against 0.8111;
Section~\ref{sec:met-limitacoes}); the default regressor is purely
synthetic, preserving the distributional part. There remains,
beyond that, the weaker risk of priors and inference policies
iteratively tuned, across versions, to performance on public
benchmarks: a mild form of community overfitting from which no
public benchmark escapes. The treatment was to declare the risk,
state conclusions at the regime level (consistency across datasets,
not individual cells), and keep the cost axis, immune to
contamination, in every practical conclusion
(Section~\ref{sec:met-limitacoes}).

\textbf{Declared asymmetries in the complementary experiments.} In
E1, NGBoost and the QRF run in documented default configuration,
without specific tuning, and the QRF receives median imputation
(fitted on the pool, applied to the test) on the three datasets with
missing values; the licensed interpretation is ``strong default'',
not ``classical tuned to the limit''
(Section~\ref{sec:destil-e1}). In the distillation students,
capacity is inherited from the benchmark's tuning, with no
re-tuning per distillation target --- a deliberate control, with
the declared caveat that between the arms both the label \emph{and}
the loss change; the asymmetry's cost was measured by the
re-tuning, registered in a dated addendum before execution, on the
20 datasets: the aggregate verdict survives the re-tuned control
(16 of 20, median 0.124 against the inherited 0.192, a declared
drop), leaving per-dataset retention dependent on the tuning budget
(Section~\ref{sec:destil-retune}). Related design-level caveats are
consolidated where they are declared: the non-nested selection
design --- a property of the benchmark (RQ1) that RQ3 inherits
through the students' hyperparameters --- with its optimism
quarantined to internal roles, in
Section~\ref{sec:met-protocolo}; the single distributional teacher
(TabPFN-2.5) and the regression-only scope, in
Chapter~\ref{chap:destilacao}'s limitations and synthesis.

\textbf{Dataset-identity errors.} Two labeling errors inherited
from the project's creation (\texttt{diamonds}/\texttt{kin8nm};
``superconduct''/\texttt{online\_news}) were discovered during the
study and corrected with re-derivation of all affected numbers;
every number in this dissertation postdates the corrections
(Section~\ref{sec:met-datasets}; record in
Appendix~\ref{app:suplementares}). The residual threat is the one
corrections do not remove: the possibility of analogous errors not
yet discovered; the mitigation is structural (data hashes,
versioned configurations, numbers traceable to artifacts).

\textbf{Multiplicity across experiments.} No correction is applied
\emph{across} experiments; under Holm correction on the
confirmatory family, the primary conclusions survive and the
supplementary instruments do not --- and RQ3's headline carries the
sequential correction of its two-stage design ($p\approx0.019$;
Section~\ref{sec:destil-desenho}).

\textbf{Undercoverage of the pure distilled student.} The
pre-registered coverage clause (PICP80 within $\pm5$~p.p.\ of
nominal) is met on only 10 of the 20 datasets --- the answer to RQ3
is qualified by this: the calibration gain is comparative, and the
closing of coverage comes from the conformal step, piloted on 5
datasets (Section~\ref{sec:destil-dist}).

\subsection{External validity}

\textbf{Small $N$ per task and the arithmetic floors.} With 10/3/5
datasets per task, the tests' power is low by construction: in
regression, the smallest two-sided Wilcoxon p-value is 0.0625
(significance at 5\% is mathematically impossible, before any
correction) and multiclass ($N{=}3$) is treated as descriptive
throughout the text. The declared mitigations (equivalences
supported by the Bayesian layer, no conclusion resting on
significance where there is no power) limit the damage but do not
eliminate it: the per-task conclusions have the precision $N$
allows, no more (Section~\ref{sec:met-limitacoes}).

\textbf{The double exclusion of \texttt{helena}.} The two
foundation models' multiclass ranks are averages over only 2
datasets, a fact repeated in every affected table. The alternative
(removing the dataset) would hide exactly the kind of architectural
restriction a practical decision study should expose; the price is
that the 2026 generation's multiclass lead rests on the smallest
sample base of the whole benchmark.

\textbf{TabFM at a pinned version, with an unreviewed technical
report.} All results refer to version v1.0.1 (July 2026, pinned
commit); the official technical report \citep{kong2026} appeared in
late September 2026 and, as of this writing, has not undergone peer
review \citep{tabfm2026}. The conclusions about the model are about
\emph{that} version; and version sensitivity is itself future work
(Section~\ref{sec:conc-futuros}).

\textbf{Delimited regime.} Eighteen OpenML datasets (plus the
fifteen of the CTR23 extension in the distillation), a
100{,}000-sample cap and a single 16~GB GPU delimit the regime
studied to small-and-medium tabular data on commonly sized
hardware. Four eligible CTR23 datasets stayed outside the
distillation population; the sensitivity with them, measured on
five data splits, keeps the verdict (17/24, $p=0.032$; conservative
cut 16/22, $p=0.026$; Chapter~\ref{chap:destilacao}). The orders of
magnitude between families tend to transfer; the absolute values do
not.

\textbf{The field moves in cycles of months.} TabPFN-2.5 in
November 2025, Pocket~FM in May 2026, TabFM in June, the partial
collision of \citet{jeong2026} in October; the difference between
the @11 and @14 columns of the LODO evaluation is the internal
demonstration of how fast conclusions date. For this reason every
novelty claim and every gap characterization carries a verification
date, from this study's dated searches, July to November 2026
(complete record in Appendix~\ref{app:suplementares}), and must be
read as ``open as far as we checked, at the declared dates'', never
as timeless certainty (Chapter~\ref{chap:related}).

\textbf{The binarized target in the regression core.} The core's
\texttt{california\_\allowbreak housing} is the binarized-target
variant of the \emph{tabular-benchmark} --- degenerate for
distributional claims. Every headline distillation result carries a
sensitivity analysis with the dataset excluded: the sign test, E3,
the five-system Friedman, and the ablations survive; the one-sided
Wilcoxon of the retention moves from $0.045$ to $0.067$ and stops
rejecting at 5\%, a degradation declared with the same prominence
as the result (Chapter~\ref{chap:destilacao}).

\textbf{Teacher licenses.} The weights of TabPFN-2.5 and TabFM are
non-commercial, and TabPFN's license further forbids distilling
from its outputs a ``competitive'' model; this study's distilled
student operates in a research regime, and commercial deployment
would require its own licensing
(Section~\ref{sec:destil-limitacoes};
Appendix~\ref{app:suplementares}).

\subsection{Construct validity}

\textbf{The \texttt{adult} latency proxy.} Latency enters the
analyses (including E2's weighted score) through the standardized
measurement in a single regime (amortized batch, each model on its
natural device), with \texttt{adult} as the reference. The bias
direction was declared in the registration and its magnitude was
later measured at the benchmark's size extremes: the rankings among
the 14 models are stable across datasets (Spearman 0.78--0.95) and
E2's crossover is identical under the three proxies
(Sections~\ref{sec:met-limitacoes}
and~\ref{sec:regret-ponderado}). What remains unmeasured is
hardware sensitivity: all numbers come from a single declared
machine.

\textbf{The grid CRPS estimator.} The 19-quantile grid
overestimates the exact CRPS by a nearly constant ${\sim}$4--5\%,
validated against the Gaussian closed form: internal comparisons
and retention ratios are preserved, but this dissertation's
absolute CRPS values should not be compared to other works'
(Section~\ref{sec:destil-metricas}).

\textbf{Metrics compress utilities.} The threshold-free primary
metric and E2's score (two criteria, linear weights) are
deliberately simple constructs of utility functions that, in
practice, include training cost, robustness, interpretability, and
regulatory constraints with non-linear weights. The framework keeps
the axes separate instead of aggregating them prematurely, but any
scalar recommendation inherits that compression
(Sections~\ref{sec:met-limitacoes}
and~\ref{sec:limitacoes-framework}).

\section{Future work}
\label{sec:conc-futuros}

Only directions with direct grounding in the chapters' evidence are
listed (extensions this study leaves ready, not speculations).

\begin{itemize}
  \item \textbf{Newer teachers: version sensitivity of the
    distillation.} The study pinned TabPFN v2.5 for comparability
    with the benchmark; TabPFN-3 and 3.5
    \citep{tabpfn3-2026,tabpfn35tr2026} already exist, and the
    partial collision of \citet{jeong2026} (point gains with newer
    teachers, under a different protocol) suggests the
    teacher--student frontier is version-sensitive. Repeating the
    distributional protocol with the next generation's teachers
    would directly measure whether the 12--19\% retention grows
    with the teacher's strength --- and whether the point negative
    resists.
  \item \textbf{Conformal calibration as complement and
    competitor.} The post-hoc conformal route on the teacher
    itself \citep{park2026} delivers coverage guarantees without a
    fast student --- at the price of the foundation model's
    latency. Incorporating it into the frontier (and applying it on
    the distilled student beyond the five-dataset pilot of
    Appendix~\ref{sec:app-pilot-conformal}) would test the
    calibration advantage observed in
    Section~\ref{sec:destil-oof} against a method with a formal
    guarantee.
  \item \textbf{Quantization as the frontier's fourth arm.}
    Attention quantization \citep{kubler2026} is the emerging
    serving lever we did not measure; the three-strategy frontier
    of Section~\ref{sec:destil-fronteira} is built to receive it
    as a fourth arm on the same hold-outs.
  \item \textbf{Multi-criteria regret with elicited preferences.}
    E2 swept linear weights over two criteria, and the latency
    proxy's robustness to the measurement dataset has been
    verified (Section~\ref{sec:regret-ponderado}); the natural
    extension is to incorporate training cost, robustness, and
    interpretability into the score and elicit utility functions
    from real practitioners --- testing whether the inversion
    point (10--20\% of latency weight, by median regret) moves with
    additional criteria.
  \item \textbf{KANs under dedicated tuning.} The negative result
    of Chapter~\ref{chap:benchmark} uses the grid's uniform
    protocol --- the fair reading for the benchmark's question,
    but not the last word on the family: a dedicated study, with
    its own search budget, the original optimizer, and basis
    variants, would definitively separate method from recipe ---
    and would say whether the bottom third is a ceiling or a
    protocol artifact.
  \item \textbf{A living update of the benchmark.} The resumable
    infrastructure, with pinned versions and seeds, is ready for
    periodic re-execution as the 2027 generation arrives;
    integration with a community-maintained living benchmark, in
    the style of TabArena \citep{tabarena2025}, is the path for
    this study's protocol depth to survive its dating --- the main
    external threat declared above.
\end{itemize}

\vspace{1em}

In synthesis, the dissertation delivers three pieces: a dated
benchmark of the 2026 generation, with the cost axes measured under
the same conditions as accuracy; a constraint framework evaluated
out of sample in its learned component, with the cost-driven
inversion point located; and a rule, tested also in its failures,
for deciding when to distill a foundation model's predictive
distribution into a student that serves at about a hundred
microseconds per row. All three depend on the declared model
versions, datasets, and hardware, and must be re-evaluated when
they change; the public protocol, with versioned seeds, artifacts,
and errata, exists so that this re-evaluation, and any correction
it brings, costs little to whoever comes next.
